\section*{Chapitre \ref{ch:g2:shapes} --- Figures et solides}

\begin{solution}{exo:g2:shapes:1}
$5$ côtés, $5$ sommets~: un pentagone. $3$ sommets (et $3$
côtés)~: un triangle. Rond~: un cercle.
\end{solution}

\begin{solution}{exo:g2:shapes:2}
\omterm{def:g2:shapes:sort}{Polygones}~: carré, triangle, hexagone (côtés droits). Pas un
\omterm{def:g2:shapes:sort}{polygone}~: le cercle (bord courbé).
\end{solution}

\begin{solution}{exo:g2:shapes:3}
Angles droits, par exemple, aux coins du tableau, d'un livre, d'une
fenêtre~; un angle non droit au couvercle ouvert d'une boîte ou
aux aiguilles d'une horloge à $2$~heures.
\end{solution}

\begin{solution}{exo:g2:shapes:4}
Carré~: $4$ angles droits. Rectangle~: $4$. Un triangle peut en
avoir au plus $1$ --- deux feraient que les deux autres côtés ne
se rencontrent jamais.
\end{solution}

\begin{solution}{exo:g2:shapes:5}
Copies sur quadrillage faites en comptant les cases entre les
sommets~; la comparaison avec le modèle vérifie le travail.
\end{solution}

\begin{solution}{exo:g2:shapes:6}
Dé~: \omterm{def:g2:shapes:solids}{cube}. Boîte de conserve~: \omterm{def:g2:shapes:solids}{cylindre}. Ballon de football~:
boule. Boîte à chaussures~: pavé.
\end{solution}

\begin{solution}{exo:g2:shapes:7}
Un \omterm{def:g2:shapes:solids}{cube} a $6$ \omterm{def:g2:shapes:solids}{faces}, chacune un carré. Un pavé a aussi $6$ \omterm{def:g2:shapes:solids}{faces}
(des rectangles).
\end{solution}

\begin{solution}{exo:g2:shapes:8}
La boule roule (ronde partout) et le \omterm{def:g2:shapes:solids}{cylindre} roule (sur son côté
courbé). Le \omterm{def:g2:shapes:solids}{cube} et le pavé ne roulent pas~: ils n'ont que des
\omterm{def:g2:shapes:solids}{faces} plates et s'arrêtent dessus.
\end{solution}

\begin{solution}{exo:g2:shapes:9}
$6$ côtés~: un hexagone, avec $6$ sommets.
\end{solution}

\begin{solution}{exo:g2:shapes:10}
Le contour de la maison (mur carré avec un toit triangulaire, la
base du toit posée sur le haut du carré) a $5$ côtés et $5$
sommets --- c'est un pentagone (bas, deux murs, deux pentes du toit~;
le haut du mur est caché sous le toit).
\end{solution}
